\scanpage{111}
\mspage{345}

\deleted{and we in fact then have
\[
g\rho(g)\rho^2(g)=g^3=\rho^{-3}=-1
\]}
\inserted{$=\begin{pmatrix}0&\unsure{-1|+1}\\1&-1\end{pmatrix}$}

\inserted{Remarks} \deleted{\unsure{N.B.}} 1°) Obviously $-\rho$ is also a solution of (\unsure{87})
(just like $\rho^{-1}$) \deleted{\unsure{and $\rho$ thus} $(-\rho)\rho(g)\rho^2(g)=1$}
(just like $-\rho^{-1}$), but it falls
respectively under cases \deleted{\illegible{1}} 1°) \deleted{\unsure{and under case}}
\deleted{\unsure{A, resp.}}

2°) To \deleted{\illegible{1}}\inserted{determine} the solutions in $\mathrm{Gl}(2,\widehat{\mathbb Z})$
of $g\rho(g)\rho^2(g)=1$, it suffices to choose for
\unsure{$\forall\,\ell$} such a solution in $\mathrm{Gl}(2,\widehat{\mathbb Z})$, using
the \unsure{result|results} (86); we shall note that
the ``cases'' obtained 1°) to 3°) can vary from one
$\ell$ to another. \deleted{Same remark for
the solutions of the equation $g\rho(g)\rho^2(g)=\unsure{-1|1}$.}
We should keep in mind that cases 1°), 2°), 3°)
correspond \unsure{resp.\textsuperscript{t}} to families of
solutions depending on 2, 1 and 0 parameters.

\unsure{3°)} I do not know whether a $\xi\in\mathrm{Sl}(2,\widehat{\mathbb Z})^{\unsure{\prime|*}}$, \inserted{$\xi=\begin{pmatrix}A&B\\C&D\end{pmatrix}$}, \inserted{which satisfies}
$g\rho(g)\rho^2(g)=1$ admits a parametric
representation of the form (58)
by $\lambda_0\rho(\lambda_0)^{-1}$ or $\lambda'_0(-l_1^{-1})\rho(\lambda'_0)^{-1}$
\inserted{$=\lambda_0(-\unsure{l_1^{-1}|l_0^{-1}})\rho(\lambda_0)^{-1}$ with $\lambda_0\unsure{=\lambda'_0\tau_\infty}$}
(it suffices
to look separately in the $\mathrm{Sl}(2,\mathbb Z_\ell)'$).
I do not know even if we assume $g\equiv\mathrm{id}\ (2)$,
(which is a condition on the component
in $\mathrm{Sl}(2,\mathbb Z_2)'$ only). We see in

\margin{\uline{N.B.} $g=l_0^{-1}$ is \uline{not}
a solution of this \unsure{the equation|equation}, on the other hand
$g=-l_0^{-1}$ is\ldots}

\scanpage{112}

\unsure{any|any} case from formulas (78) that \unsure{\uline{if}} $\xi$
\deleted{condition}\inserted{is of} the form $\lambda_0\rho(\lambda_0)^{-1}$
\inserted{(if $\xi\rho(\xi)\rho^2(\xi)=1$, of course)},
\deleted{\unsure{then indeed has case A)}} \deleted{\illegible{1}}
\inserted{we will have} $C+D-B=1$ --- i.e. \deleted{\illegible{1}}\inserted{we are}
\inserted{in the} ``general'' case with two parameters.
\deleted{I presume
that if we take
\[
g=\lambda_0(-l_0^{-1})\rho(\lambda_0)^{-1},
\]
we will be in case B), with $c+d-b=-1$
i.e. case 2°) again with two parameters. I
\unsure{dispense with} doing the calculation here.}

\margin{OK, cf. (\unsure{99}) below.}

(Without nevertheless being able to exclude the case
\inserted{where we should at the same time}
$A+B=0$, \deleted{which in fact means that
\[
\xi=\begin{pmatrix}A&B\\C&D\end{pmatrix}
=-\rho^{-1}=\begin{pmatrix}-1&1\\-1&0\end{pmatrix}
\]}
i.e. the case with \unsure{one} parameter
\[
\xi=\begin{pmatrix}-u&u\\C&-u'-C\end{pmatrix}
\]
with $u$ a primitive \unsure{third} root of 1,
$u'=u^{-1}=u^2$ the other primitive root.

\scanpage{113}
\mspage{346}

(VIII) --- \uline{Calculations in} $\mathrm{Sl}(2,\widehat{\mathbb Z})$ : \uline{cocycles under} $\sigma_\infty$.

Let us take again a $g\in\begin{pmatrix}a&b\\c&d\end{pmatrix}\in\mathrm{Sl}(2,\widehat{\mathbb Z})$,
and compute $\sigma_\infty g\sigma_\infty^{-1}$, we find, thanks to
\[
\sigma_\infty=\begin{pmatrix}0&-1\\1&0\end{pmatrix}\qquad (48)
\]
\deleted{\illegible{1}}
the formulas
\[
(90)\quad\left\{
\begin{aligned}
\sigma_\infty g\sigma_\infty^{-1}
&=\begin{pmatrix}d&-c\\-b&a\end{pmatrix}={}^t(g^{-1})
&&\text{the contragredient of $g$ !}\\
&\text{i.e.}\\
\sigma_\infty g^{-1}\sigma_\infty^{\unsure{-1}}
&={}^t g&&\text{thus}\\
[\sigma_\infty,g^{-1}]&=({}^t g)(g)
=\begin{pmatrix}a^2+b^2&ac+bd\\ac+bd&\unsure{c^2+d^2|c^2-d^2}\end{pmatrix}
\end{aligned}\right.
\]
\inserted{$=\sigma_\infty g^{-1}\unsure{\sigma_\infty|\sigma_\infty^{-1}}g$}

\margin{\uline{N.B.} we have
$\sigma_\infty^{-1}=\begin{pmatrix}0&1\\-1&0\end{pmatrix}=-\sigma_\infty$
and \deleted{\illegible{1}} $\sigma_\infty^2\ne1$ in $\mathrm{Sl}(2,\widehat{\mathbb Z})$ !}

\margin{\uline{N.B.} If $g\in\mathrm{Gl}(2,A)$,
$A$ any ring
(for example $A=\mathbb Z_\ell$), we have
$\sigma_\infty g\sigma_\infty^{-1}=\begin{pmatrix}d&-c\\-b&a\end{pmatrix}
=(\det g){}^t(g^{-1})$}

thus, taking the inverse of both sides in
the last formula
\[
(91)\quad \underbrace{g^{-1}\sigma_\infty(g)}_{\inserted{({}^t g\cdot g)^{-1}}}
=\begin{pmatrix}c^2+d^2&-(ac+bd)\\-(ac+bd)&a^2+b^2\end{pmatrix}.
\]

On the other hand, let
\deleted{$h_0=\begin{pmatrix}\unsure{r}&s\\t&u\end{pmatrix}$}
$h_0\in\mathrm{Sl}(2,\widehat{\mathbb Z})$
we have
\[
h_0\sigma_\infty(h_0)=h_0({}^t h_0^{-1})
\]
thus
\[
(*)\qquad h_0\sigma_\infty(h_0)=1\quad\Longleftrightarrow\quad h_0={}^t h_0.
\]

Finally, let
\[
(91)\qquad \xi=\begin{pmatrix}A&B\\C&D\end{pmatrix}
\]
and consider the equation (68)
\[
(**)\qquad \sigma_\infty(\xi)=l_1^\gamma\xi^{-1}l_0^\gamma
\]
which is also equivalent to (*), with $h_0$ given by (\unsure{65|66})
\deleted{$h_0=l_0^{-\unsure{\gamma|\delta}}\xi$ \illegible{4}}, $l_0,l_1$ by (44).

\scanpage{114}

\[
(92)\quad\begin{aligned}
h_0=l_0^{-\gamma}\xi
&=\begin{pmatrix}1&+2\gamma\\0&1\end{pmatrix}
\begin{pmatrix}A&B\\C&D\end{pmatrix}\\
&=\begin{pmatrix}A+2\gamma C&B+2\gamma D\\C&D\end{pmatrix}
\end{aligned}
\]
thus ``the equation in $\xi$'' (**) (written in the
quotient $\mathrm{Sl}(2,\widehat{\mathbb Z})'$ of $\widehat{\mathfrak S}_{0,3}^{+}$, is written $B+2\gamma D=C$
i.e.
\[
(93)\qquad \underbrace{C-B}_{\inserted{1-D}}=2\gamma D
\qquad\text{i.e.}\quad\boxed{pD=1}\ \inserted{\text{\unsure{in case|one case}}}
\]
\deleted{\illegible{2}}

\margin{on \unsure{recalling} that $2\gamma+1=p$}

(We should keep in mind that the equation (**) \uline{in}
$\mathrm{Sl}(2,\widehat{\mathbb Z})'$ means $h_0\sigma_\infty(h_0)=\pm1$ i.e. $h_0=\pm{}^t h_0$,
\deleted{thus here} but the sign $-$ would imply that we have
$D=-D$ i.e. $2D=0$ so $D=0$, which is impossible
if we assume that $\xi=\begin{pmatrix}A&B\\C&D\end{pmatrix}$ comes
from $\widehat\Pi_{0,3}$ so is in $\Pi'(2)$ i.e.
$\equiv\deleted{\unsure{\pm}}1\ (2)$ (so that $D\equiv1\ (2)$).
Thus we must indeed have (93).

Let us now take the parametric
representation \inserted{(70)} of $\xi$ (taking account of
the equation $\xi\rho(\xi)\rho^2(\xi)=1$ \deleted{\unsure{$\xi$}} (64))
\[
(94)\qquad\begin{aligned}
\xi&=\lambda\rho(\lambda)^{-1}
&&\inserted{(\lambda\in\mathrm{Sl}(2,\widehat{\mathbb Z}))}\\
&\text{or}\quad
\xi=\lambda'(\unsure{-l_1^{-1}})\rho(\lambda')^{-1}
=\lambda\rho(\lambda)^{-1}
&&\inserted{\left(\begin{gathered}\lambda=\lambda'\tau_\infty\\\lambda'\in\mathrm{Sl}(2,\widehat{\mathbb Z})\end{gathered}\right)}
\end{aligned}
\]
\deleted{\unsure{$\lambda,\lambda'\in\mathrm{Sl}$}}
with $\lambda\in\mathrm{Gl}(2,\widehat{\mathbb Z})^{\unsure{*|\circ}}$ \inserted{\deleted{\unsure{case I}} $\det\lambda=\pm1$},
(in fact $\lambda\equiv\deleted{\unsure{\pm}}1\ (2)$ in the
case that interests us, but no matter\ldots)

\deleted{A) Case $\xi=\lambda\rho(\lambda)^{-1}=[\lambda,\rho]$ with $\lambda=\begin{pmatrix}a&b\\c&d\end{pmatrix}$,}

\scanpage{115}
\mspage{347}

then formula (78) gives $A,B,C,D$ as
quadratic functions in $a,b,c,d$, \deleted{we will have}
\inserted{(paying Attention \unsure{to the} sign $\delta=-1$ in case 2°)}
\deleted{
\[
(95)\quad\left\{
\begin{aligned}
\unsure{B-C|B+C}&=\unsure{\delta(c^2+d^2+cd-\delta)}=D-1\\
D&=\delta(c^2+d^2+cd)
\end{aligned}\right.
\]
\unsure{$D=$}}
thus (93) is written
\deleted{$1-D=2\gamma D$ \unsure{or} i.e.}
\[
(96)\qquad\boxed{(2\gamma+1)D=1\qquad\text{i.e.}\quad pD=1}
\qquad D=\delta(c^2+d^2+cd),
\]
i.e. the multiplier $p$ is obtained as a function
of $\xi\in\widehat\Pi_{0,3}$ by taking the image
of $\xi$ in $\mathrm{Sl}(2,\widehat{\mathbb Z})'$, in the form $\begin{pmatrix}A&B\\C&D\end{pmatrix}$
(there is a sign ambiguity, removed by the
choice of $\lambda$ such that $\xi=\lambda\rho(\lambda)^{-1}$, the image
of $\lambda$ in $\mathrm{Gl}(2,\widehat{\mathbb Z})'$ \unsure{not|in} determining a
matrix (still denoted $\lambda$) of $\mathrm{Gl}(2,\widehat{\mathbb Z})$ except up to a
\uline{sign}, but $\lambda\rho(\lambda)^{-1}$ no longer depends on the
choice of this \uline{sign}\ldots), by the formula
\[
(97)\qquad p=\frac1D\qquad (D=\unsure{\delta}(c^2+d^2+cd))
\]

Let us now note that on reduction modulo 3,
the quadratic form $c^2+d^2+cd=(c+d)^2-cd$
takes only the values 0 and 1 (clear if $c=0$,
if $d=0$, \deleted{or $c+d=0$ then $f(c,d)=-cd=c^2$}
--- so we
find a perfect square in these \deleted{three} cases, otherwise
$c,d=\pm1$, and $f(c,d)=1+1\pm1\equiv1$ or $0$ \deleted{\illegible{1}} (3). Now here
we must find a unit so the value 0 excluded,
so

\scanpage{116}

\margin{\uline{N.B.} We have $D\equiv p\equiv\delta\ (3)$ (98) knowledge of $D$ (or $p$, or $\delta$) modulo $3$ \deleted{\unsure{already|thus} tells} \unsure{tells} us --- which case we have in the parametric \unsure{representation} of $\xi$. Useless now to treat \unsure{these} cases separately --- since we have $\xi=\lambda\rho(\lambda)^{-1}$, $\lambda=\lambda'\tau_\infty$, $\det\lambda=-1$.}

\[
p\equiv\delta\quad(3)
\]
\inserted{(where $\delta=\det\lambda\in\{\pm1\}$, $\xi=\lambda\rho(\lambda)^{-1}$)}
\deleted{\illegible{4} $\xi=\lambda\rho(\lambda)^{-1}$, $\det\lambda=\delta\in\{\pm1\}$, \illegible{3}}
i.e.
\[
\left\{\begin{aligned}
\text{Case }\xi&=\lambda\rho(\lambda)^{-1}\quad(\lambda\in\mathrm{Sl}(2,\widehat{\mathbb Z}))
\Longleftrightarrow p\equiv1\ (3)\quad\deleted{\unsure{g=\lambda_0\rho(\lambda_0)^{-1}}}\Longleftrightarrow\gamma\equiv0\ (3)\\
\text{Case }\xi&=\lambda'l_1^{-1}\unsure{\rho(\lambda)^{-1}|\rho(\lambda')^{-1}}\quad(\lambda'\in\mathrm{Sl}(2,\widehat{\mathbb Z}))
\Longleftrightarrow p\equiv-1\ (3)\Longleftrightarrow\gamma\equiv-1\ (3).
\end{aligned}\right.
\]

B) Case $\xi=\lambda'l_0^{-1}\rho(\lambda')^{-1}=(\lambda'l_0^{-1})(\rho\lambda'^{-1}\rho^{-1})$.

Using the second formula (77) we find
\[
\xi=\begin{pmatrix}a&b\\c&d\end{pmatrix}
\begin{pmatrix}1&2\\0&1\end{pmatrix}
\begin{pmatrix}a-c&c\\a+b-c-d&c+d\end{pmatrix}
=\begin{pmatrix}A&B\\C&D\end{pmatrix}
\]
with
\[
(99)\quad
\left\{\begin{aligned}
A&=3a^2+b^2+3ab-3ac-2ad-bc-bd
& B&=3ac+2ad+bc+bd\\
C&=-(3c^2+3cd+\deleted{\unsure{2}}d^2)+3ac+ad+2bc+\deleted{\unsure{2}}bd
& D&=3c^2+3cd+d^2
\end{aligned}\right.,
\]
\margin{\uline{N.B.} $C+D-B=\unsure{-1|-\delta}$.}

and formula (93) (which expresses ``the equation in $\xi$'' (**), or the one in $h_0$ (*)) is written, \deleted{\unsure{thus again in the form}} taking account of
\[
B-C=-1+D,
\]
in the form \deleted{\unsure{$2\gamma$}}
\[
2\gamma D=-1-D\qquad\text{i.e.}
\]
\[
(100)\quad \boxed{(2\gamma+1)D=-1\qquad\text{i.e.}\quad pD=-1},
\]
\margin{as before we cannot have $h_0=-{}^t h_0$, for we cannot have $D=0$ if we \unsure{assume} \unsure{that} $\xi\equiv1\ (2)$, \unsure{hence} $D\equiv1\ (2)$.}

so here the multiplier $p$ is expressed as a function of the image of $\xi$ in $\mathrm{Sl}(2,\widehat{\mathbb Z})'$, put in the form $\begin{pmatrix}A&B\\C&D\end{pmatrix}$ where $A,B,C,D\in\widehat{\mathbb Z}$ --- and on removing the sign ambiguity \unsure{as before} --- by the formula
\[
(101)\quad p=-\frac1D\qquad(D=\unsure{3c^2+3cd+d^2|3(c^2+3cd+d^2)}).
\]